\section{Introduction}
\label{sec:introduction}

Gaussian variational inference replaces an intractable target distribution
\(\pi(x)\propto e^{-V(x)}\) on \(\R^d\) by the minimizer of
\[
  \calE(m,C)
  :=\KL\!\left(\N(m,C)\,\middle\|\,\pi\right),
  \qquad (m,C)\in\R^d\times\SPD(d).
\]
The approximation is computationally useful only if the optimization respects
the strong coupling between the mean and covariance.  Euclidean updates do not:
their behavior changes under an invertible affine reparameterization of the
target.  The Fisher--Rao natural gradient is invariant under precisely these
changes of coordinates and is therefore a canonical geometry for the full
Gaussian family \citep{rao1945information,amari1998natural,amari2016information}.
This paper develops a global-to-local rate theory and identifies several rates
that this geometry can and cannot deliver.

The main conclusion is a global-to-local rate picture with three distinct
bottlenecks.  First, a small initial covariance suppresses the mean motion and
creates a covariance burn-in.  Second, after the covariance has reached the
target curvature scale, global localization is controlled by an intrinsic
condition number obtained by whitening with the optimizing covariance.  Third,
inside a neighborhood of the optimizer, the rate is governed by a score-
operator spectral gap that measures non-Gaussian coupling between the mean and
covariance modes.  These mechanisms are logically independent.  A transport
warm start removes the first, but neither the global lower construction nor the
local spectral obstruction disappears.

\subsection{The three bottlenecks}
\label{sec:introduction-bottlenecks}

Assume initially that
\begin{equation}
  \alpha\Id\preceq\nabla^2V(x)\preceq\beta\Id,
  \qquad \kappa:=\beta/\alpha.
  \label{eq:introduction-curvature}
\end{equation}
Let \(a_\star=(m_\star,C_\star)\) be the optimizing Gaussian.  In the
optimizer-whitened coordinates
\[
  z=C_\star^{-1/2}(x-m_\star),\qquad
  \widehat V(z)=V(m_\star+C_\star^{1/2}z),
\]
write \(\alphastar,\betastar\) for the best curvature bounds of
\(\widehat V\), and \(\kappastar=\betastar/\alphastar\).  These quantities are
affine invariant, satisfy \(1\leq\kappastar\leq\kappa^2\), and equal one for
every Gaussian target, even if its original precision is arbitrarily
ill-conditioned.

The covariance bootstrap is visible directly in precision coordinates.  If
\(\widehat C_0=C_\star^{-1/2}C_0C_\star^{-1/2}\) and
\(\lambda_{0,\star}=\lambda_{\min}(\widehat C_0)\), then the Fisher--Rao flow
satisfies
\begin{equation}
  \calE(a_t)-\calE(a_\star)
  \leq
  \Delta_0
  \left[1+\betastar\lambda_{0,\star}(e^t-1)\right]^{-1/\kappastar}.
  \label{eq:introduction-global-rate}
\end{equation}
The first logarithm in the resulting hitting time depends on
\(\lambda_{0,\star}\); the subsequent localization time is proportional to
\(\kappastar\).  A two-dimensional logarithmic spiral forces
\(\Omega(\kappastar)\) continuous time under an admissible anisotropic
initialization.  For fixed-step Riemannian and KL schemes, a one-dimensional
bump train, run with \(h=\gamma/\kappa\) for
\(0<\gamma\leq1/2\), forces \(\Omega(\kappa/h)\) iterations for a
constant-factor reduction even after covariance burn-in.  The two
constructions therefore isolate distinct continuous-time and fixed-step
obstructions.

At the optimizer, strong log-concavity alone does not reduce the problem to a
Gaussian quadratic.  The first- and second-Hermite components of the whitened
Hessian form a self-adjoint linearized generator.  We prove the spectral
sandwich
\begin{equation}
  \max\!\left\{\alphastar,\frac1{4+\Gamma}\right\}
  \leq\gstar\leq\Lstar
  \leq\min\!\left\{\betastar,\frac{4+\Gamma}{2}\right\},
  \qquad
  \Gamma=O(1+\sqrt d\log\kappastar).
  \label{eq:introduction-local-gap}
\end{equation}
The lower bound is dimension dependent, and this is not merely a proof
artifact in growing dimension.  A smooth convex-ridge family with
\(d=\Theta(\kappastar^2)\) satisfies
\(\gstar=\Theta(\kappastar^{-1/2})\) and a squared first-Hermite coupling of
order \(\kappastar^{1/2}\).  It saturates the
\(\betastar-1\) branch of \(\Gamma\).  It does not settle whether a stronger
bound holds in fixed dimension or whether the
\(\sqrt d\log\kappastar\) branch is sharp.

\subsection{Contributions}
\label{sec:introduction-contributions}

The deterministic theory gives both upper and lower bounds.  The continuous
estimate \eqref{eq:introduction-global-rate} is affine invariant and separates
initial covariance from intrinsic conditioning.  The two deterministic
discretizations admit two-sided covariance bootstraps and have the common
complexity
\[
  O\!\left(
    h^{-1}\logp\frac1{\beta\lambda_{\min}(C_0)}
    +\frac{\kappa}{h}\logp\frac{\Delta_0}{\delta}
  \right),
\]
with a lower bound matching the \(\kappa/h\) factor for constant-factor
reduction.  A Bures--Wasserstein warm start or one pointwise
gradient--Hessian quadratic rescue can remove the covariance term, but not the
other two stages.

The local theory keeps the nonlinear Gaussian covariance dynamics exactly and
perturbs only the non-Gaussian part.  This Gaussian-core decomposition gives an
explicit energy-sublevel region for the flow, constructive local regions for
both maps, exact continuous and KL invariant regions on Gaussian targets, and
an online residual gate for entry.  For general discrete targets, the radius
depends on a supremum of second derivatives of the update map over a compact
hull.  The result is constructive but is not advertised as a closed-form
practical radius.

The stochastic theory uses a joint Price/Hessian oracle and a
sticking-the-landing mean estimator \citep{roeder2017sticking}.  It supplies
pathwise covariance safety after the quadratic rescue under the stated
stepsizes and global constant-step bounds with explicit residual levels.  Under
Hessian-Lipschitzness it also gives a floor-free \(O(1/N)\) decreasing-step
guarantee.  Under the same assumption, the local result is deliberately
separated into a killed-process mean-square contraction and a finite-horizon
containment probability.  Under the universal spectral certificate its
iteration dependence is quadratic, rather than linear, in \(4+\Gamma\).
Infinite-horizon containment is not claimed.

Two extensions mark the scope of the theory.  First, for \(d\geq2\) we
classify the complete three-parameter family of affine-invariant Riemannian
metrics on Gaussians.  On a Gaussian target the parameters precondition the
mean, trace, and traceless covariance modes separately; Fisher--Rao is the
unique choice, up to scale, that balances all three.  This modal statement does
not extend every sharp rate theorem to every metric.  Second, a smooth
non-log-concave double well produces a Riemannian covariance cascade and a pole
in the KL resolvent.  A covariance-constrained KL step nevertheless has an
\(O(1/N)\) stationarity guarantee in a corrected split displacement.  General
Wasserstein--Fisher--Rao dynamics are left to a separate paper.

\subsection{Relation to prior work}
\label{sec:introduction-related}

Natural gradients and information geometry provide the classical background
for the Fisher--Rao metric \citep{rao1945information,amari1998natural,amari2016information}.
Gaussian variational inference and stochastic variational inference are
reviewed in \citet{blei2017variational} and
\citet{JMLR:v14:hoffman13a}.  The present analysis differs from a generic
natural-gradient convergence argument by exploiting the exact Gaussian
covariance equation and by retaining affine-invariant, optimizer-whitened
quantities throughout.

Wasserstein Gaussian variational methods offer a complementary geometry.
\citet{lambert2022variational} study variational inference through Wasserstein
gradient flows, while \citet{diao2023forward} analyze a Bures--Wasserstein
forward--backward scheme.  We use Bures--Wasserstein only as a warm start and
as a nonconvex comparison; it is not blended with Fisher--Rao into a general
WFR algorithm here.  Affine-invariant metrics on positive-definite matrices
were classified in several forms by
\citet{thanwerdas2019affine,thanwerdas2023n}.  Our classification includes the
mean block of the Gaussian family and, for \(d\geq2\), uses it to expose the
exact trace and traceless modal rates.

The Price/Hessian stochastic construction is related to the sticking-the-
landing principle of \citet{roeder2017sticking}.  Recent natural-gradient VI
work also studies nonconjugate models and constrained covariance schemes
\citep{sun2025natural}.  Our stochastic results focus on the joint geometry of
the mean and full covariance, distinguish constant-step residual bounds from
decreasing-step convergence, and keep local containment separate from local
contraction.

\subsection{Organization and notation}
\label{sec:introduction-organization}

\Cref{sec:geometry} introduces the objective, the affine metric family, and the
algorithms.  \Cref{sec:global} develops the deterministic global theory and its
lower bounds.  \Cref{sec:local} gives the local spectral and nonlinear theory.
\Cref{sec:stochastic} treats Price/Hessian algorithms.
\Cref{sec:nonconvex} identifies the non-log-concave boundary, and
\Cref{sec:numerics} presents focused numerical illustrations.

We write \(d\) for the ambient dimension, \(\norm{\cdot}_{\op}\) and
\(\norm{\cdot}_{\Frob}\) for the operator and Frobenius norms, and
\(\Delta(a)=\calE(a)-\calE(a_\star)\).  The symbol \(\kappa\) always denotes
the original-coordinate curvature ratio in
\eqref{eq:introduction-curvature}; \(\kappastar\) denotes its
optimizer-whitened counterpart.  Constants hidden by \(O(\cdot)\) are
numerical unless their dependence is stated, and
\(\logp x:=\max\{\log x,0\}\).
